% !TEX program = lualatex
\documentclass[12pt,a4paper]{article}

\usepackage{luatexja}
\usepackage{amsmath,amssymb,amsthm}
\usepackage{geometry}
\usepackage{listings}
\usepackage{xcolor}
\usepackage{booktabs}
\usepackage{hyperref}
\usepackage{enumitem}
\usepackage{algorithm}
\usepackage{algpseudocode}

\theoremstyle{definition}
\newtheorem{defn}{Definition}[section]
\newtheorem{remark}{Remark}[section]

\algrenewcommand\algorithmicrequire{\textbf{Input:}}
\algrenewcommand\algorithmicensure{\textbf{Output:}}
\newcommand{\algorithmicbreak}{\textbf{break}}
\newcommand{\Break}{\State\algorithmicbreak}
\floatname{algorithm}{Algorithm}

\geometry{margin=25mm}

\hypersetup{
  colorlinks=true,
  linkcolor=blue!60!black,
  urlcolor=blue!60!black
}

\lstset{
  basicstyle=\ttfamily\small,
  keywordstyle=\color{blue!70!black}\bfseries,
  commentstyle=\color{green!50!black},
  stringstyle=\color{red!60!black},
  frame=single,
  breaklines=true,
  breakatwhitespace=false,
  showstringspaces=false,
  tabsize=2,
  numbers=left,
  numberstyle=\tiny\color{gray},
  numbersep=6pt,
  xleftmargin=12pt,
  keepspaces=true
}

\title{\textbf{Apple-Peel Unfolding of Four-Dimensional Regular Polytopes}\\[0.4em]
  \large Algorithm Design and Implementation Notes}
\author{Takashi Yoshino}
\date{\today}

\begin{document}
\maketitle
\tableofcontents
\bigskip


% ============================================================
\section{Background and Motivation}
% ============================================================

\subsection{What Is Apple-Peel Unfolding?}

\textbf{Apple-Peel Unfolding} is a method for selecting the faces (or cells) of a
polyhedron or polytope one by one in a spiral order---analogous to peeling an
apple in a single continuous strip---so as to produce an unfolding (net).
The method was formalised for three-dimensional polyhedra by
Yoshino, Chaidee, and Sawatdithep in \texttt{applePeelingV01.pdf}.

The goal of this project is to extend the method to
\textbf{four-dimensional regular polytopes}.

\subsection{Target Polytopes}

The six regular convex 4-polytopes are studied (Table~\ref{tab:polytopes}).

\begin{table}[h]
\centering
\caption{Basic data of the regular 4-polytopes}
\label{tab:polytopes}
\small
\begin{tabular}{llrrrr}
\toprule
Name & $\{p,q,r\}$ & Cells & Cell type & Vertices & Data file \\
\midrule
5-cell (pentachoron)            & $\{3,3,3\}$ &   5 & tetrahedron      &   5 & \texttt{f5.m}  \\
8-cell / hypercube (tesseract)  & $\{4,3,3\}$ &   8 & cube             &  16 & \texttt{f8.m}  \\
16-cell (hexadecachoron)        & $\{3,3,4\}$ &  16 & tetrahedron      &   8 & \texttt{f16.m} \\
24-cell (icositetrachoron)      & $\{3,4,3\}$ &  24 & octahedron       &  24 & \texttt{f24.m} \\
120-cell (hecatonicosachoron)   & $\{5,3,3\}$ & 120 & dodecahedron     & 600 & \texttt{f120.m}\\
600-cell (hexacosichoron)       & $\{3,3,5\}$ & 600 & tetrahedron      & 120 & \texttt{f600.m}\\
\bottomrule
\end{tabular}
\end{table}

\subsection{References}

\begin{itemize}
  \item Yoshino, Chaidee, Sawatdithep:
        \textit{Apple-Peel Unfolding of Polyhedra}
        (\texttt{applePeelingV01.pdf}) --- original 3D formulation.
  \item Kaino (2019): prior work constructing 4D nets of the 24-cell,
        120-cell, and 600-cell geometrically by orthogonal projection.
\end{itemize}


% ============================================================
\section{Overview of the 3D Algorithm}
% ============================================================

\subsection{Algorithm 1 (applePeelingV01)}

\begin{algorithm}[h]
\caption{Apple-Peel Unfolding (3D version)}
\label{alg:3d}
\begin{algorithmic}[1]
\Require
  Vertex coordinates $\mathit{vers}$ (list of 3D points),
  face list $\mathit{faces}$ (each entry: list of vertex indices forming one face),
  adjacency list $\mathit{nbrs}$ ($\mathit{nbrs}[i]$: indices of faces adjacent to face $i$),
  starting face $F_1$, second face $F_2$
\Ensure
  Face selection order $\mathit{order}$ starting from $(F_1, F_2)$;
  success flag $\mathit{success}$

\State $\mathit{vers} \gets \mathit{vers} - \bar{\mathit{vers}}$
  \Comment{Translate so that the centroid of all vertices is at the origin}
\State $\hat{c}_1 \gets$ centroid of face $F_1$, normalised to a unit vector
\State $\mathit{vers} \gets \mathrm{Rotate}(\mathit{vers},\;\hat{c}_1 \to +\hat{z})$
  \Comment{Rotate so that $\hat{c}_1$ aligns with the $+z$ axis}
\State $\mathit{order} \gets [F_1,\; F_2]$
\State $\hat{c}_i \gets$ centroid of face $i$ \textbf{for all} $i$
  \Comment{Recompute centroids after rotation}
\State $\mathit{last} \gets F_2$
  \Comment{Index of the most recently selected face}
\While{there exists an unselected adjacent face}
  \State $L \gets \bigl\{j \in \mathit{nbrs}[\mathit{last}]
           \mid \hat{c}_{\mathit{last},y}\,\hat{c}_{j,x}
                - \hat{c}_{\mathit{last},x}\,\hat{c}_{j,y} \geq 0\bigr\}$
    \Comment{Left half-space candidates ($x,y$: first two coordinates of $\hat{c}$)}
  \If{$L \neq \emptyset$}
    \State $\mathit{next} \gets \arg\max_{j \in L}\; \hat{c}_{j,z}$
      \Comment{Among left candidates, choose the one with the largest $z$ coordinate}
  \Else
    \State $\mathit{next} \gets \arg\min_{j \in \mathit{nbrs}[\mathit{last}]}\; \hat{c}_{j,z}$
      \Comment{Fallback: choose the candidate with the smallest $z$}
  \EndIf
  \State $\mathit{order}.\mathrm{append}(\mathit{next})$
  \State $\mathit{nbrs}[i] \gets \mathit{nbrs}[i] \setminus \{\mathit{next}\}$
    \hspace{1.5em}\textbf{for all} $i$
    \Comment{Remove $\mathit{next}$ from every adjacency list}
  \State $\mathit{last} \gets \mathit{next}$
\EndWhile
\State \Return $\bigl(\mathit{order},\;|\mathit{order}| = |\mathit{faces}|\bigr)$
\end{algorithmic}
\end{algorithm}

\subsection{Peelability Classification}

\begin{itemize}
  \item \textbf{Perfect}: every $(F_1, F_2)$ pair yields a complete unfolding.
  \item \textbf{Possible}: some pairs succeed and some fail.
  \item \textbf{Impossible}: no $(F_1, F_2)$ pair yields a complete unfolding.
\end{itemize}

\noindent
For the 13 Archimedean solids the reported counts are:
Perfect~4, Possible~3, Impossible~6.

\subsection{Mathematica Implementation \texttt{peeling3Df}}

The implementation resides in function \texttt{peeling3Df} inside
\texttt{221026ArchimedeanSolidsPt02.nb}.
Its signature is:

\begin{lstlisting}
peeling3Df[origVers_, faces_, top_]
\end{lstlisting}

\noindent
Here \texttt{origVers} is the list of vertex coordinates, \texttt{faces} is
the list of faces (each face given as a list of vertex indices), and
\texttt{top} is the index of the starting face $F_1$.
Adjacent faces are detected with a \texttt{RotateLeft}-based edge-pair trick.
The fallback (when only right-side candidates exist) takes
\texttt{neighbors[[lastChosen]][[1]]} (the first element of the list),
which is a simplification of the ``min-$z$'' rule in the paper.

The return value is, for each $F_2$, a triple
\texttt{\{rotatedVers, faceOrder, successBool\}}.


% ============================================================
\section{Extension to Four Dimensions}
% ============================================================

\subsection{Coordinate System and Peeling Axis}

The four-dimensional space is coordinatised as $O\text{-}xyzw$.
The \textbf{$w$-axis} serves as the peeling axis, playing the role of the
$z$-axis in the 3D algorithm.

\subsection{4D Data Structure}

Data are stored in \texttt{\_4DData/f[n].m}.
Loading the file returns a six-element list:

\begin{center}
\texttt{\{vers, edgs, neis, faces, fneis, cells\}}
\end{center}

\noindent
The meaning of each element is given in Table~\ref{tab:data}.

\begin{table}[h]
\centering
\caption{Data structure of \texttt{f[n].m}}
\label{tab:data}
\begin{tabular}{lll}
\toprule
Variable & Shape & Contents \\
\midrule
\texttt{vers}  & $(\text{\#vertices}) \times 4$ & 4D vertex coordinates \\
\texttt{edgs}  & $(\text{\#edges}) \times 2$    & Pairs of vertex indices forming each edge \\
\texttt{neis}  & $(\text{\#vertices}) \times *$ & Adjacent vertex indices for each vertex \\
\texttt{faces} & $(\text{\#faces}) \times *$    & Vertex indices forming each 2D face \\
\texttt{fneis} & $(\text{\#faces}) \times *$    & Adjacent face indices for each face \\
\texttt{cells} & $(\text{\#cells}) \times *$    & \textbf{Face indices forming each 3D cell} \\
\bottomrule
\end{tabular}
\end{table}

\noindent
\textbf{Important}: \texttt{cells[[i]]} contains \textbf{face indices}
(not vertex indices).
The vertex set of cell $i$ is therefore

\begin{equation}
  V_i = \bigcup_{f \in \texttt{cells}[[i]]} \texttt{faces}[[f]].
\end{equation}

Example (5-cell, $n=5$):
\begin{lstlisting}
cells = {{1,2,4,7}, {1,3,5,8}, {2,3,6,9}, {4,5,6,10}, {7,8,9,10}}
(* each cell is a tetrahedron consisting of 4 triangular faces *)
\end{lstlisting}

\subsection{Analogy between 3D and 4D}

\begin{center}
\begin{tabular}{ll}
\toprule
3D & 4D \\
\midrule
Face $F$ of a polyhedron & Cell $C$ of a 4-polytope (a 3D polyhedron) \\
Peeling axis $z$ & Peeling axis $w$ \\
Face centroid $\hat{c} \in \mathbb{R}^3$ & Cell centroid $\hat{c} \in \mathbb{R}^4$ \\
Rotate $F_1$ toward $+z$ & Rotate $C_1$ toward $+w$ \\
Left half-space $(c_k \times \hat{z}) \cdot c' \geq 0$ &
  Left half-space $c_y c'_x - c_x c'_y \geq 0$ \\
Adjacency: $\geq 2$ shared vertices & Adjacency: $\geq 1$ shared face \\
\bottomrule
\end{tabular}
\end{center}

\subsection{Derivation of the Left Half-Space Condition}

Expanding the 3D condition $(\hat{c}_k \times \hat{z}) \cdot \hat{c}' \geq 0$
with $\hat{z} = (0,0,1)$ gives:
\[
  (c_{k,y},\, -c_{k,x},\, 0) \cdot c' = c_{k,y}\,c'_x - c_{k,x}\,c'_y \geq 0.
\]
Here $c_{k,x}, c_{k,y}$ are the $x$- and $y$-components of the centroid
$\hat{c}_k$ of the current face, and $c'_x, c'_y$ are the $x$- and
$y$-components of the centroid of the candidate face.
Because this expression depends only on the $xy$-components of the two
centroids, it can be \textbf{applied directly in four dimensions}:
\[
  \boxed{c_y\,c'_x - c_x\,c'_y \geq 0}.
\]
This is equivalent to a rotation in the $xy$-plane (about the $zw$-plane)
and is the most compact 4D extension.

\subsection{Implementation of the 4D Rotation}

To rotate the vertices so that the centroid $\hat{c}_1$ of cell $C_1$ points
in the $+w$ direction, we use Mathematica's \texttt{RotationMatrix}:

\begin{lstlisting}
angle4D = ArcCos[Clip[cc1 . wHat / Norm[cc1], {-1, 1}]];
RotationMatrix[angle4D, {cc1/Norm[cc1], wHat}] . # & /@ localVers
\end{lstlisting}

\noindent
Here \texttt{cc1} is the centroid of $C_1$ and
\texttt{wHat} $= (0,0,0,1)$ is the unit vector along the $+w$ axis.
\texttt{RotationMatrix[}\textit{angle}\texttt{, \{u, v\}]} generates a rotation
matrix in the plane spanned by unit vectors $u$ and $v$,
which works in any dimension and therefore applies directly to 4D.


% ============================================================
\section{Re-examination of the Left Condition and Rigorous 4D Extension}
\label{sec:left}
% ============================================================

This section formalises the ``choose the leftmost cell'' condition of the
3D Apple-Peel algorithm and derives its correct extension to four dimensions.

% ----------------------------------------------------------
\subsection{Geometric Meaning of the Left Condition in 3D}
% ----------------------------------------------------------

After aligning the peeling axis with the $z$-axis, each step selects the
neighbour of the current face whose centroid lies farthest to the ``left''.
``Left'' is defined by the $xy$-plane angle:

\begin{equation}
  c_{k,y}\, c'_x - c_{k,x}\, c'_y \;\ge\; 0,
  \label{eq:3dleft}
\end{equation}
where $c_{k,x}, c_{k,y}$ are the $x$- and $y$-components of $\mathbf{c}_k$
(centroid of the current face at step $k$), and $c'_x, c'_y$ are
the $x$- and $y$-components of $\mathbf{c}'$ (centroid of the candidate face).

\paragraph{Interpretation via the Darboux frame.}
Introduce a local frame (Darboux frame) at each step:
\begin{itemize}
  \item Normal direction:
    $\hat{n}_k = \mathbf{c}_k / |\mathbf{c}_k|$
    (position on the unit sphere; $\mathbf{c}_k$ is the centroid of the
     current face, treated as a point on a sphere centred at the origin).
  \item Forward direction:
    $\hat{f}_k = (\mathbf{c}_{k-1} \times \mathbf{c}_k) /
                  |\mathbf{c}_{k-1} \times \mathbf{c}_k|$,
    where $\mathbf{c}_{k-1}$ is the centroid of the \emph{previous} face.
    This cross product lies in the tangent plane at $\hat{n}_k$ and points
    in the ``left'' direction orthogonal to the path from $\mathbf{c}_{k-1}$
    to $\mathbf{c}_k$.
  \item Left direction:
    $\hat{L}_k = \hat{n}_k \times \hat{f}_k$
    (90° to the left of the forward direction within the tangent plane).
\end{itemize}

The candidate face $j$ with centroid $\mathbf{c}_j$ lies in the
\emph{left half-space} if and only if $\hat{L}_k \cdot \mathbf{c}_j \ge 0$.
This is exactly equivalent to the determinant condition
\begin{equation}
  \det\!\bigl(\mathbf{c}_{k-1},\, \mathbf{c}_k,\, \mathbf{c}_j\bigr) \;\ge\; 0,
  \label{eq:det3d}
\end{equation}
because
\[
  \hat{L}_k \cdot \mathbf{c}_j
  = \frac{\det(\mathbf{c}_{k-1}, \mathbf{c}_k, \mathbf{c}_j)}
         {|\mathbf{c}_k|\,|\hat{d}_k^\perp|},
\]
where $\hat{d}_k^\perp$ is the normalised component of
$\mathbf{c}_{k-1}$ perpendicular to $\mathbf{c}_k$
(the scalar factor is positive), so the signs of
\eqref{eq:3dleft} and \eqref{eq:det3d} always agree.

\begin{remark}
  In three dimensions, $\mathbf{c}_{k-1} \times \mathbf{c}_k$ yields a
  unique one-dimensional ``left'' direction.
  This is why ``left'' is well-defined in 3D.
\end{remark}

% ----------------------------------------------------------
\subsection{Difficulty in 4D: Dimension of the Orthogonal Complement}
% ----------------------------------------------------------

In $\mathbb{R}^4$, the orthogonal complement of the 2D subspace spanned by
$\mathbf{c}_{k-1}$ and $\mathbf{c}_k$ is \textbf{2-dimensional}
(in $\mathbb{R}^3$ it is 1-dimensional).
Hence no unique ``left vector'' analogous to the cross product exists,
and the notion of ``left'' is not canonically determined.

To circumvent this, we adopt the following approach.
In the 4D Apple-Peel, we project each cell centroid onto the
$xyz$-hyperplane (discarding the $w$-component) to obtain a 3D vector
$\tilde{\mathbf{c}}_k = (c_{k,x},\, c_{k,y},\, c_{k,z}) \in \mathbb{R}^3$,
and apply the determinant condition in this projected space.

% ----------------------------------------------------------
\subsection{Rigorous 4D Extension: Determinant Condition Using the Previous Step}
% ----------------------------------------------------------

\begin{defn}[4D Apple-Peel Left Condition]
  \label{defn:4dleft}
  At step $k \ge 2$, let
  $\tilde{\mathbf{c}}_k \in \mathbb{R}^3$ be the $xyz$-projection of the
  centroid of the \emph{current} cell,
  $\tilde{\mathbf{c}}_{k-1} \in \mathbb{R}^3$ be the $xyz$-projection of the
  centroid of the \emph{previous} cell, and
  $\tilde{\mathbf{c}}_j \in \mathbb{R}^3$ be the $xyz$-projection of the
  centroid of \emph{candidate cell} $j$.
  Cell $j$ is said to lie in the \textbf{left half-space} if
  \begin{equation}
    \det\!\bigl(\tilde{\mathbf{c}}_{k-1},\, \tilde{\mathbf{c}}_k,\, \tilde{\mathbf{c}}_j\bigr)
    \;\ge\; 0.
    \label{eq:4dleft}
  \end{equation}
\end{defn}

This condition is equivalent to the left direction of the Darboux frame
restricted to the $xyz$-hyperplane.

\paragraph{Ambiguity at the initial step.}
At step $k = 1$ (the very first selection after $C_2$ is fixed),
$\tilde{\mathbf{c}}_{k-1}$ does not yet exist, so \eqref{eq:4dleft} cannot
be applied.
The current implementation (\texttt{peeling4Df.m}) uses at $k = 1$ the
2-component condition
$c_{k,y} c'_x - c_{k,x} c'_y \ge 0$,
which corresponds to taking $(0,0,1)$ as the implicit reference vector
$\tilde{\mathbf{c}}_{k-1}$.
The choice of initial condition may affect the results.

% ----------------------------------------------------------
\subsection{Three Selection Rules}
% ----------------------------------------------------------

Given that multiple candidates satisfy the left half-space condition,
the following three rules define how to select the next cell.
Viewing the $xyz$-projected centroids as points on the sphere $S^2$,
the geometric meaning of each rule becomes clear.
Let $\hat{n} = \tilde{\mathbf{c}}_k / |\tilde{\mathbf{c}}_k|$
be the unit vector of the current position,
$\hat{f}$ be the forward direction of the Darboux frame,
and $\varphi_j$ be the signed angle between the projection of
$\tilde{\mathbf{c}}_j$ onto the tangent plane and $\hat{f}$
(positive $=$ to the left).

\begin{defn}[Selection Rules]
  \label{defn:rules}
  Let $\mathcal{L}_k$ denote the set of left half-space candidates.
  \begin{itemize}
    \item \textbf{Rule~1 (max angle):}
      $\displaystyle j^* = \arg\max_{j \in \mathcal{L}_k} \varphi_j.$
      Select the candidate making the sharpest left turn.
      A large $\varphi_j$ (close to $90°$) corresponds to a tight spiral.
    \item \textbf{Rule~2 (min angle / loxodrome):}
      $\displaystyle j^* = \arg\min_{j \in \mathcal{L}_k} \varphi_j.$
      Select the candidate closest to the forward direction while still
      turning left.
      This approximates a \textbf{loxodrome (rhumb line)} on the sphere:
      the path maintains a constant small left bearing relative to the
      direction of travel.
    \item \textbf{Rule~3 (max $w$ / latitude-first):}
      $\displaystyle j^* = \arg\max_{j \in \mathcal{L}_k} \hat{c}_{j,w},$
      where $\hat{c}_{j,w}$ is the $w$-coordinate of the centroid of
      candidate cell $j$.
      This selects the cell at the highest ``latitude'' (greatest $w$),
      corresponding to a \textbf{latitude-by-latitude scan}: complete one
      latitude band before descending to the next.
  \end{itemize}
\end{defn}

% ----------------------------------------------------------
\subsection{Geometric Comparison of Rule~2 (Loxodrome) and Rule~3 (Latitude-First)}
% ----------------------------------------------------------

\paragraph{Comparison on the 3D sphere.}
In three dimensions, Rule~2 traces a loxodrome (equiangular spiral) on
the sphere, while Rule~3 traces a path that completes each latitude circle
before descending.
The two are generally different:
\begin{itemize}
  \item A \textbf{loxodrome} maintains a constant angle with every meridian
        and spirals infinitely toward the pole (approximated on a finite polyhedron).
  \item The \textbf{latitude-first path (Rule~3)} is a discrete analogue that
        traverses cells at a fixed latitude before moving to a lower latitude
        (smaller $w$).
\end{itemize}

\paragraph{Difference on polytopes.}
On a discrete polytope only finitely many cells are available at each step,
so neither rule exactly reproduces its continuous counterpart.
Nevertheless, the difference in selection criteria can lead to different
success rates.

% ----------------------------------------------------------
\subsection{Correspondence with Kaino (2019)}
% ----------------------------------------------------------

Kaino (2019) constructs 4D nets by scanning the faces of a sphere-inscribed
regular polytope latitude by latitude.
This structure corresponds to \textbf{Rule~3 (max $w$)} in our framework:
start from the high-latitude cells near $C_1$ (top) and descend to lower
latitudes (smaller $w$) step by step.

The current \texttt{peeling4Df.m} selects cells in decreasing order of
$\hat{c}_{j,w}$ among left candidates, matching Rule~3 conceptually.
The one difference is that the implementation uses the 2-component left
condition \eqref{eq:3dleft} rather than the rigorous Definition~\ref{defn:4dleft}.

% ----------------------------------------------------------
\subsection{Proposed Two-Stage Selection}
% ----------------------------------------------------------

As a result of the above discussion, we propose the following two-stage
strategy:

\begin{table}[h]
\centering
\caption{Proposed two-stage selection strategy}
\label{tab:twostage}
\begin{tabular}{cll}
\toprule
Stage & Condition & Action \\
\midrule
1st & Try Rule~3 (max $w$) &
  Select the left candidate with the largest $w$ (current implementation) \\
2nd & Fall back to Rule~2 (loxodrome) &
  If Rule~3 gets stuck, re-select using the min-$\varphi$ criterion \\
\bottomrule
\end{tabular}
\end{table}

Rationale:
\begin{itemize}
  \item Rule~3 is simple to define and consistent with Kaino's approach.
  \item Rule~2 is theoretically interesting (discrete approximation of a
        spherical loxodrome), but Rule~3 more intuitively reproduces a
        spiral peel on a polytope.
  \item A staged approach---try the simpler rule first, then apply the more
        refined one on failure---is also pragmatically efficient.
\end{itemize}

Systematically comparing the success rates of Rules 1, 2, and 3 on small
polytopes (5-cell, 8-cell, etc.) to identify the best strategy is a future
task (see \S\ref{sec:future}).


% ============================================================
\section{Implementation of \texttt{peeling4Df}}
% ============================================================

\subsection{Function Signature}

\begin{lstlisting}
peeling4Df[origVers_, faces_, cells_, top_]
\end{lstlisting}

\noindent Arguments:
\begin{itemize}
  \item \texttt{origVers}: list of 4D vertex coordinates
        (one 4-element vector per vertex).
  \item \texttt{faces}: list of 2D faces; \texttt{faces[[f]]} is a list of
        vertex indices forming face $f$.
  \item \texttt{cells}: list of 3D cells; \texttt{cells[[i]]} is a list of
        \emph{face indices} forming cell $i$.
  \item \texttt{top}: index of the starting cell $C_1$.
\end{itemize}

Return value: for each $C_2$ candidate (each neighbour of $C_1$), a triple
\texttt{\{rotatedVers, selectionOrder, successBool\}}, where
\texttt{rotatedVers} is the vertex array after the $+w$ rotation,
\texttt{selectionOrder} is the list of cell indices in the order they were
selected, and \texttt{successBool} is \texttt{True} iff all cells were
selected.

\subsection{Algorithm}

\begin{algorithm}[h]
\caption{peeling4Df (4D greedy version)}
\label{alg:4d}
\begin{algorithmic}[1]
\Require 4D vertex coordinates $\mathit{vers}$, face list $\mathit{faces}$,
         cell list $\mathit{cells}$, starting cell $\mathit{top}$
\Ensure For each $C_2$: $(v',\; \mathit{order},\; \mathit{success})$

\Function{CellVerts}{$i$}
  \Return $\bigcup_{f \in \mathit{cells}[i]} \mathit{faces}[f]$
    \Comment{Vertex index set of cell $i$}
\EndFunction
\Statex
\State $v \gets \mathit{vers} - \bar{\mathit{vers}}$
  \Comment{$\bar{\mathit{vers}}$: mean of all vertex coords; translate to centroid at origin}
\State $\hat{c}_1 \gets \overline{v[\textsc{CellVerts}(\mathit{top})]}$
  \Comment{Centroid of starting cell $C_1$ in the translated frame}
\State $v \gets \mathrm{Rotate4D}(v,\; \hat{c}_1 \to +\hat{w})$
  \Comment{4D rotation aligning $\hat{c}_1$ with the $+w$ axis}
\For{$i = 1$ \textbf{to} $|\mathit{cells}|$}
  \State $\hat{c}_i \gets \overline{v[\textsc{CellVerts}(i)]}$
    \Comment{Centroid of cell $i$ after rotation ($\hat{c}_i \in \mathbb{R}^4$)}
\EndFor
\For{$i = 1$ \textbf{to} $|\mathit{cells}|$}
  \State $\mathit{adj}[i] \gets
    \{j \neq i : \mathit{cells}[i] \cap \mathit{cells}[j] \neq \emptyset\}$
    \Comment{$\mathit{adj}[i]$: indices of cells sharing at least one face with cell $i$}
\EndFor
\Statex
\For{each $C_2 \in \mathit{adj}[\mathit{top}]$}
  \State $\mathit{order} \gets [\mathit{top},\; C_2]$
    \Comment{$\mathit{order}$: selected cell sequence, initialised with $C_1$ and $C_2$}
  \State $\mathit{nbrs} \gets$ copy of $\mathit{adj}$, with $\mathit{top}$ and $C_2$ removed
    \Comment{$\mathit{nbrs}[i]$: unselected neighbours of cell $i$}
  \State $\mathit{last} \gets C_2$
    \Comment{$\mathit{last}$: most recently selected cell}
  \While{$\mathit{nbrs}[\mathit{last}] \neq \emptyset$}
    \State $L \gets \bigl\{j \in \mathit{nbrs}[\mathit{last}] \mid
             \hat{c}_{\mathit{last},y}\,\hat{c}_{j,x}
             - \hat{c}_{\mathit{last},x}\,\hat{c}_{j,y} \geq 0\bigr\}$
      \Comment{$L$: left half-space candidates; subscripts $x,y$ denote coordinate components}
    \If{$L \neq \emptyset$}
      \State $\mathit{next} \gets \arg\max_{j \in L}\; \hat{c}_{j,w}$
        \Comment{$\hat{c}_{j,w}$: $w$-coordinate of centroid of cell $j$}
    \Else
      \State $\mathit{next} \gets \arg\min_{j \in \mathit{nbrs}[\mathit{last}]}\; \hat{c}_{j,w}$
        \Comment{Fallback: among right-side candidates choose the one with min $w$}
    \EndIf
    \State $\mathit{order}.\mathrm{append}(\mathit{next})$
    \State $\mathit{nbrs}[i] \gets \mathit{nbrs}[i] \setminus \{\mathit{next}\}$
      \hspace{1em}\textbf{for all} $i$
    \State $\mathit{last} \gets \mathit{next}$
  \EndWhile
  \State \textbf{yield} $\bigl(v,\; \mathit{order},\; |\mathit{order}| = |\mathit{cells}|\bigr)$
\EndFor
\end{algorithmic}
\end{algorithm}

\subsection{Key Differences from the 3D Implementation}

\begin{itemize}
  \item Cell vertices are retrieved as
        \texttt{Union @@ faces[[ cells[[i]] ]]},
        because \texttt{cells[[i]]} stores face indices, not vertex indices.
  \item Adjacency is defined by sharing $\geq 1$ face (not $\geq 2$ vertices
        as in 3D), because two adjacent 3D cells share a 2D face.
  \item The fallback is \textbf{exactly} ``min $w$ among right-side
        candidates'' (the 3D implementation used the simpler
        \texttt{neighbors[[lastChosen]][[1]]}).
  \item The 4D rotation uses \texttt{RotationMatrix[angle, \{u, v\}]}
        with 4D vectors.
\end{itemize}


% ============================================================
\section{File Structure and Usage}
% ============================================================

\subsection{File List}

\begin{table}[h]
\centering
\caption{Main files}
\begin{tabular}{ll}
\toprule
File & Description \\
\midrule
\texttt{peeling4Df.m}       & Core algorithm (use via \texttt{Get}) \\
\texttt{run4DPeeling.m}     & Script to run and classify all 6 polytopes \\
\texttt{4DPeelingAll.nb}    & Notebook combining the above two files \\
\texttt{\_4DData/f[n].m}    & Data for each regular polytope ($n=5,8,16,24,120,600$) \\
\texttt{applePeelingV01.pdf}& Paper describing the 3D algorithm \\
\bottomrule
\end{tabular}
\end{table}

\subsection{How to Run}

\paragraph{Method 1 --- Via Notebook}
\begin{lstlisting}
(* Open 4DPeelingAll.nb in Mathematica and choose
   Evaluation -> Evaluate Notebook *)
\end{lstlisting}

\paragraph{Method 2 --- Script execution}
\begin{lstlisting}
SetDirectory["/path/to/260324Peeling4D"];
Get["peeling4Df.m"];
Get["run4DPeeling.m"];
\end{lstlisting}

\paragraph{Method 3 --- Direct function call}
\begin{lstlisting}
(* Load the 5-cell data *)
{vers, edgs, neis, faces, fneis, cells} =
    Get["/path/to/_4DData/f5.m"];

(* Run peeling with starting cell 1 *)
result = peeling4Df[N[vers], faces, cells, 1];

(* result[[k]] = {rotatedVers, order, successBool} for the k-th C2 *)
result[[1, 3]]   (* True or False *)
result[[1, 2]]   (* cell selection order *)
\end{lstlisting}

\subsection{Checking the Classification}

\begin{lstlisting}
(* Collect results for all (C1, C2) pairs *)
ans = Table[peeling4Df[N[vers], faces, cells, top],
            {top, Length[cells]}];

allBools  = Flatten[Map[Last, ans, {2}]];
trueCount = Count[allBools, True];

classification = Which[
  Count[allBools, False] == 0, "Perfect",
  trueCount == 0,              "Impossible",
  True,                        "Possible"
];
\end{lstlisting}


% ============================================================
\section{Computational Results}
% ============================================================

\subsection{3D Regular Polyhedra}
\label{ssec:res3dplatonic}

Table~\ref{tab:res3dplatonic} summarises the success rates for the five
regular polyhedra under the current algorithm (global $c_1$ reference,
$\varepsilon = 10^{-10}$, no single-candidate bypass).
RS (Spiral rule) minimises $\det(\mathbf{c}_1,\mathbf{c}_k,\mathbf{c}_j)$
(most negative = sharpest CW turn; fallback uses Darboux frame $\min\varphi$);
RZ (Zonal rule) maximises $z$ with $\min\det$ as secondary criterion.
``w'' = with fallback; ``n'' = no fallback.

\begin{table}[h]
\centering
\caption{Success rates for 3D regular polyhedra}
\label{tab:res3dplatonic}
\begin{tabular}{llrcccc}
\toprule
Polyhedron & $\{p,q\}$ & Pairs & RS(w) & RS(n) & RZ(w) & RZ(n) \\
\midrule
Tetrahedron  & $\{3,3\}$ & 12 & 100\% & 100\% & 100\% & 100\% \\
Cube         & $\{4,3\}$ & 24 & 100\% & 100\% & 100\% & 100\% \\
Octahedron   & $\{3,4\}$ & 24 & 100\% & 100\% & 100\% & 100\% \\
Dodecahedron & $\{5,3\}$ & 60 & 100\% & 100\% & 100\% & 100\% \\
Icosahedron  & $\{3,5\}$ & 60 & 100\% & 100\% & 100\% & 100\% \\
\bottomrule
\end{tabular}
\end{table}

\noindent
All five polyhedra achieve 100\% under both rules, with and without
fallback.
This all-or-nothing pattern is consistent with the equivariance theorem:
since regular polyhedra are face-transitive, success of any $(F_1,F_2)$
pair implies success of all symmetry-related pairs.
(Under the strict variant $\varepsilon<0$, the no-fallback rates drop to
0\% for all solids except the Tetrahedron; see the detailed discussion
in the algorithm section.)

\subsection{3D Archimedean Solids}
\label{ssec:res3darchimedean}

Table~\ref{tab:res3darchimedean} shows the success rates for the thirteen
Archimedean solids computed with global $c_1$ reference.
``w'' uses the fallback branch (\texttt{run\_archi\-medean\_faceup.m});
``n'' exits immediately when no left candidate exists
(\texttt{run\_archi\-medean\_nofallback.m}).

\begin{table}[h]
\centering
\caption{Success rates (\%) for 3D Archimedean solids.
  ``w'' = with fallback; ``n'' = no fallback.}
\label{tab:res3darchimedean}
\small
\setlength{\tabcolsep}{4pt}
\begin{tabular}{llrcccc}
\toprule
& & & \multicolumn{2}{c}{RS\,(\%)} & \multicolumn{2}{c}{RZ\,(\%)} \\
\cmidrule(lr){4-5}\cmidrule(lr){6-7}
Polyhedron & v.c. & Pairs & w & n & w & n \\
\midrule
Truncated Tetrahedron       & 3.6.6       &  36 & 66.7 & 66.7 & 33.3 & 33.3 \\
Cuboctahedron               & 3.4.3.4     &  48 &  0.0 &  0.0 &  0.0 &  0.0 \\
Truncated Cube              & 3.8.8       &  72 &  0.0 &  0.0 &  0.0 &  0.0 \\
Truncated Octahedron        & 4.6.6       &  72 & 41.7 & 41.7 & 100  & 100  \\
Rhombicuboctahedron         & 3.4.4.4     &  96 &  0.0 &  0.0 &  0.0 &  0.0 \\
Truncated Cuboctahedron     & 4.6.8       & 144 &  0.0 &  0.0 & 100  & 100  \\
Snub Cube                   & 3.3.3.3.4   & 120 & 20.0 &  0.0 & 40.0 & 20.0 \\
Icosidodecahedron           & 3.5.3.5     & 120 &  0.0 &  0.0 &  0.0 &  0.0 \\
Truncated Dodecahedron      & 3.10.10     & 180 &  0.0 &  0.0 &  0.0 &  0.0 \\
Truncated Icosahedron       & 5.6.6       & 180 &  0.0 &  0.0 & 100  & 100  \\
Rhombicosidodecahedron      & 3.4.5.4     & 240 &  0.0 &  0.0 &  0.0 &  0.0 \\
Truncated Icosidodecahedron & 4.6.10      & 360 &  0.0 &  0.0 & 66.7 & 66.7 \\
Snub Dodecahedron           & 3.3.3.3.5   & 300 &  0.0 &  0.0 &  0.0 &  0.0 \\
\bottomrule
\end{tabular}
\end{table}

\noindent
Key observations:
\begin{itemize}
  \item Seven solids (54\%) fail under both rules (0\%).
  \item \textbf{Three solids achieve 100\% under RZ}:
        Truncated Octahedron, Truncated Cuboctahedron, and
        Truncated Icosahedron.
        The Truncated Icosidodecahedron achieves 66.7\% under RZ
        (RS = 0\%).
  \item \textbf{RS outperforms RZ on one solid}:
        Truncated Tetrahedron (RS = 66.7\% $>$ RZ = 33.3\%).
  \item \textbf{RZ $>$ RS on four solids}:
        Truncated Octahedron (100\% vs.\ 41.7\%),
        Truncated Cuboctahedron (100\% vs.\ 0\%),
        Truncated Icosahedron (100\% vs.\ 0\%), and
        Truncated Icosidodecahedron (66.7\% vs.\ 0\%).
  \item \textbf{RZ results are invariant under the switch from
        min-$\varphi$ to min-Det} (the two tiebreaks are equivalent).
        RS results changed significantly (min-Det $\ne$ min-$\varphi$).
  \item \textbf{The fallback is essential only for the Snub Cube.}
        Removing it leaves all 12 other solids unchanged (0\% stays 0\%,
        100\% stays 100\%).
        For the Snub Cube: RS drops from 20.0\% (24/120) to 0\%,
        and RZ drops from 40.0\% (48/120) to 20.0\% (24/120).
\end{itemize}

\paragraph{Equivariance verification.}
For each solid we checked whether pairs $(F_1, F_2)$ with the same face-type
combination (polygon size of $F_1$, polygon size of $F_2$) always give the
same outcome (all succeed or all fail).
12 of the 13 solids pass this check for both RS and RZ.

The sole exception is the Snub Cube: within the $(3,3)$ face-type group
the 72 triangle-to-triangle pairs split $24/72$ rather than $0/72$ or $72/72$.
This is not a numerical inconsistency.
The Snub Cube has two symmetrically distinct types of triangles under its
chiral octahedral symmetry group $O$ (order 24):
\emph{snub triangles} (24 faces, each sharing one edge with a square)
and \emph{gyrate triangles} (8 faces, sharing no edge with a square).
Splitting by sub-orbit yields three perfectly uniform groups:
$\{\text{snub},\text{snub}\}: 24/24$,
$\{\text{snub},\text{gyrate}\}: 0/24$, and
$\{\text{gyrate},\text{snub}\}: 0/24$.

\emph{Antipodal-case bug and fix.}
Investigation of the 4 failures in the $\{4,3\}$ (square--triangle) group
revealed a bug in \texttt{p3lAlignTopToZ}: when the top-face centroid points
exactly in the $-z$ direction, the code used coordinate inversion $-\mathbf{v}$,
which in $\mathbb{R}^3$ has $\det(-I_3)=-1$ (improper rotation) and reverses
the left half-space condition.
The fix replaces it with a proper $180^\circ$ rotation
$(x,y,z)\mapsto(x,-y,-z)$ about the $x$-axis, which has $\det=+1$.
After the fix, $\{\mathrm{sq},\mathrm{snub}\}$ becomes $24/24$
and the overall RS/RZ rate rises from 36.7\% to 40.0\%.

\paragraph{Mirror symmetry verification.}
A reflection $\rho\colon(x,y,z)\mapsto(-x,y,z)$ has $\det(\rho)=-1$, so it
reverses the sign of the left half-space determinant.
Running the left-spiral algorithm on the mirror image is therefore equivalent
to running a right-spiral algorithm on the original polyhedron.

We ran RS and RZ on both the original and the mirror image of all 13
Archimedean solids.
\textbf{All 13 solids yield identical success counts for original and mirror image}
(Table~\ref{tab:mirror}).

\begin{table}[h]
\centering
\caption{Success counts: original vs.\ mirror image (all 13 Archimedean solids).}
\label{tab:mirror}
\small
\setlength{\tabcolsep}{4pt}
\begin{tabular}{llrcccc}
\toprule
Solid & v.c. & Pairs & \multicolumn{2}{c}{RS} & \multicolumn{2}{c}{RZ} \\
\cmidrule(lr){4-5}\cmidrule(lr){6-7}
 & & & Orig & Mirror & Orig & Mirror \\
\midrule
Truncated Tetrahedron       & 3.6.6       &  36 & 24 & 24 & 12 & 12 \\
Cuboctahedron               & 3.4.3.4     &  48 &  0 &  0 &  0 &  0 \\
Truncated Cube              & 3.8.8       &  72 &  0 &  0 &  0 &  0 \\
Truncated Octahedron        & 4.6.6       &  72 & 30 & 30 & 72 & 72 \\
Rhombicuboctahedron         & 3.4.4.4     &  96 &  0 &  0 &  0 &  0 \\
Truncated Cuboctahedron     & 4.6.8       & 144 &  0 &  0 &144 &144 \\
Snub Cube$^*$               & 3.3.3.3.4   & 120 & 24 & 24 & 48 & 48 \\
Icosidodecahedron           & 3.5.3.5     & 120 &  0 &  0 &  0 &  0 \\
Truncated Dodecahedron      & 3.10.10     & 180 &  0 &  0 &  0 &  0 \\
Truncated Icosahedron       & 5.6.6       & 180 &  0 &  0 &180 &180 \\
Rhombicosidodecahedron      & 3.4.5.4     & 240 &  0 &  0 &  0 &  0 \\
Truncated Icosidodecahedron & 4.6.10      & 360 &  0 &  0 &240 &240 \\
Snub Dodecahedron$^*$       & 3.3.3.3.5   & 300 &  0 &  0 &  0 &  0 \\
\bottomrule
\multicolumn{7}{l}{\small$^*$ Chiral solid.}
\end{tabular}
\end{table}

For the 11 amphichiral solids the mirror image is congruent to the original
under a proper rotation, so the equality of counts follows immediately from
equivariance.

For the two chiral solids, the Snub Dodecahedron yields 0\% in both cases
(trivial agreement).
For the \emph{Snub Cube}, both RS and RZ preserve their success counts
under mirror reflection (RS: 24/120 each; RZ: 48/120 each),
but the underlying pairs differ.
Under RZ, both original and mirror give 48/120 (40.0\%) but
the successful $(F_1,F_2)$ pairs differ: 24 pairs succeed under both spirals
($\{\mathrm{snub},\mathrm{snub}\}$), 24 succeed only under the left spiral
($\{\mathrm{sq},\mathrm{snub}\}$), and 24 succeed only under the right spiral
($\{\mathrm{gyrate},\mathrm{snub}\}$).
Under RS, both original and mirror give 24/120 (20.0\%);
only the $\{\mathrm{snub},\mathrm{snub}\}$ orbit (24/24) succeeds,
and the same orbit succeeds for both the original and its mirror.
The orbits $\{\mathrm{sq},\mathrm{snub}\}$ and $\{\mathrm{gyrate},\mathrm{snub}\}$
are exchanged by the reflection, and their equal sizes (24 pairs each) explain
why the RZ count is preserved.
The algorithm is thus sensitive to chirality at the orbit level, correctly
distinguishing the two enantiomorphs of the Snub Cube even though the total
count is the same.

\subsection{4D Regular Polytopes: Classification}
\label{ssec:res4d}

Running \texttt{peeling4Df4} (global $c_1$--$c_2$ reference) under both
the Zonal rule (RZ: max-$w$ primary, Det tie-break) and the Spiral rule
(RS: max-Det primary, max-$w$ tie-break)
over all $(C_1, C_2)$ pairs gives the classification shown in
Table~\ref{tab:results}.

\begin{table}[h]
\centering
\caption{Apple-Peel unfolding of regular 4-polytopes (\texttt{peeling4Df4}, RZ vs.\ RS).
  ``RZ'' = successful pairs under the Zonal rule; ``RS'' = successful pairs under
  the Spiral rule. True $=$ Unique for all polytopes; ``Class'' refers to RZ.}
\label{tab:results}
\small
\setlength{\tabcolsep}{3pt}
\begin{tabular}{llrrrrrl}
\toprule
Name & $\{p,q,r\}$ & Cells & Adj. & $(C_1,C_2)$ & RZ & RS & Class (RZ) \\
\midrule
5-cell (pentachoron)           & $\{3,3,3\}$ &   5 &  4 &    20   &     20 &     20 & \textbf{Perfect}    \\
8-cell (tesseract)             & $\{4,3,3\}$ &   8 &  6 &    48   &     48 &     48 & \textbf{Perfect}    \\
16-cell (hexadecachoron)       & $\{3,3,4\}$ &  16 &  4 &    64   &     25 &      0 & \textbf{Possible}   \\
24-cell (icositetrachoron)     & $\{3,4,3\}$ &  24 &  8 &   192   &    192 &    192 & \textbf{Perfect}    \\
120-cell (hecatonicosachoron)  & $\{5,3,3\}$ & 120 & 12 & 1{,}440 & 1{,}250 &      0 & \textbf{Possible}   \\
600-cell (hexacosichoron)      & $\{3,3,5\}$ & 600 &  4 & 2{,}400 &      0 &      0 & \textbf{Impossible} \\
\bottomrule
\end{tabular}
\end{table}

\noindent
\textbf{RZ vs.\ RS in 4D.}
RZ substantially outperforms RS.
RS is Impossible on the 16-cell (0/64) and the 120-cell (0/1\,440),
while RZ achieves 25/64 (39.1\%) and 1\,250/1\,440 (86.8\%), respectively.
On the 5-cell, 8-cell, and 24-cell both rules produce identical success counts.

\smallskip\noindent
\textbf{Notes.}
The 8-cell is a special case: all cell centroids lie on coordinate axes
($\pm e_i$), so $\det(c_1,c_2,c_k,c_j)\equiv 0$ and every candidate
passes the filter at every step, reducing the selection to max-$w$ alone;
RZ and RS agree trivially.
The 600-cell is Impossible under both rules.
All 2\,400 evaluated pairs terminate at one of exactly five step values
(146, 150, 276, 279, or 284), indicating structural bottlenecks
tied to the 600-cell's icosahedral symmetry.

\subsection{Comparison with Previous 4D Algorithm (\texttt{peeling4Df})}
\label{ssec:res4dcomp}

Table~\ref{tab:res4dcomp} compares unique success counts between
\texttt{peeling4Df} (xy two-component cross product, local reference)
and \texttt{peeling4Df4} (4D determinant condition, global $c_1$--$c_2$
reference).

\begin{table}[h]
\centering
\caption{Unique success counts: \texttt{peeling4Df} vs.\ \texttt{peeling4Df4} (Det tie-break)}
\label{tab:res4dcomp}
\begin{tabular}{llrrl}
\toprule
Name & $\{p,q,r\}$ & \texttt{peeling4Df} & \texttt{peeling4Df4} & Change \\
\midrule
5-cell   & $\{3,3,3\}$ &   20 &       20 & same \\
8-cell   & $\{4,3,3\}$ &   48 &       48 & same (degenerate$^\dagger$) \\
16-cell  & $\{3,3,4\}$ &   26 &       25 & $-1$ \\
24-cell  & $\{3,4,3\}$ &  119 &      192 & $+73$ (Perfect) \\
120-cell & $\{5,3,3\}$ &  357 & 1{,}250 & $+893$ \\
600-cell & $\{3,3,5\}$ & --- (incomplete) &        0 & Impossible \\
\bottomrule
\end{tabular}
\end{table}

\noindent
$^\dagger$ The 8-cell determinant degenerates to zero; both algorithms
reduce to max-$w$ selection.

The global reference with geometric tie-break substantially improves
the 24-cell ($+73$, now Perfect) and 120-cell ($+893$),
while the 16-cell is essentially unchanged ($-1$).
The 600-cell, which could not be completed by \texttt{peeling4Df} within
the allotted time, is confirmed to be Impossible under the global
condition.

\subsection{Success Rates}
\label{ssec:successrates}

\begin{table}[h]
\centering
\caption{$(C_1, C_2)$ success rates (\texttt{peeling4Df4}, RZ vs.\ RS)}
\label{tab:successrate}
\begin{tabular}{lrrl}
\toprule
Name & RZ rate & RS rate & Class (RZ) \\
\midrule
5-cell   & 100.0\% ($20/20$)     & 100.0\% ($20/20$)   & Perfect    \\
8-cell   & 100.0\% ($48/48$)     & 100.0\% ($48/48$)   & Perfect    \\
16-cell  &  39.1\% ($25/64$)     &   0.0\% ($0/64$)    & Possible   \\
24-cell  & 100.0\% ($192/192$)   & 100.0\% ($192/192$) & Perfect    \\
120-cell &  86.8\% ($1250/1440$) &   0.0\% ($0/1440$)  & Possible   \\
600-cell &   0.0\% ($0/2400$)    &   0.0\% ($0/2400$)  & Impossible \\
\bottomrule
\end{tabular}
\end{table}

\noindent
In 4D, RZ dominates RS.
RZ achieves Perfect on the 24-cell (100\%) and 86.8\% on the 120-cell,
while RS is Impossible on both the 16-cell and the 120-cell.
The 5-cell, 8-cell, and 24-cell show identical success counts for both rules.


% ============================================================
\section{3D Visualisation and Verification of Unfoldings}
% ============================================================

\subsection{4D $\to$ 3D Unfolding (\texttt{unfold3DExport.m})}

Given a cell selection order, the implementation in
\texttt{unfold3DExport.m} places each cell (a 3D polyhedron) in 3D space
sequentially to form a net.
The procedure is as follows.

\begin{enumerate}[label=\textbf{Step \arabic*.}]
  \item \textbf{SVD projection}:
        Project the 4D vertices of each cell onto the best-fitting 3D
        hyperplane using singular value decomposition, obtaining local
        3D coordinates that preserve distances.

  \item \textbf{Procrustes alignment}:
        For each adjacent pair of cells, find the rotation $R$ and
        translation $\boldsymbol{t}$ (no scaling) that best aligns the
        shared face vertices in 3D.

  \item \textbf{Reflection correction}:
        After Procrustes alignment, if the newly placed cell and its
        predecessor lie on the \emph{same} side of the shared face, reflect
        the new cell through that face to place it on the ``outside''.

  \item \textbf{STL export}:
        Triangulate the unfolded cells and export the result as an STL file
        using \texttt{Graphics3D} + \texttt{Export}.
\end{enumerate}

\subsection{Overlap Check (Separating Axis Theorem)}

Whether distinct cells overlap in the 3D unfolding is tested for every pair
using the Separating Axis Theorem (SAT).

\begin{itemize}
  \item Candidate axes: face normals of each cell and cross products of
        edge vectors.
  \item AABB (axis-aligned bounding box) pre-filter for speed.
  \item Tolerance $\varepsilon = 10^{-6}$: contact along a shared face is
        not counted as overlap.
\end{itemize}

\begin{table}[h]
\centering
\caption{Number of generated STL files (\texttt{STLs/} folder)}
\label{tab:stlcount}
\begin{tabular}{lr}
\toprule
Polytope & STL count \\
\midrule
5-cell   &   20 \\
8-cell   &   48 \\
16-cell  &   26 \\
24-cell  &  119 \\
120-cell &  357 \\
600-cell &  --- (incomplete) \\
\midrule
Total    &  570 \\
\bottomrule
\end{tabular}
\end{table}

\noindent
The STL files include nets that contain overlaps.
Counting the valid nets (no overlap) is a future task
(\S\ref{sec:future}).

\subsection{4D Net Visualisation (\texttt{unfold4D.m})}
\label{ssec:unfold4D}

The script \texttt{unfold4D.m} reads the peeling results stored in
\texttt{ans4DGlobal.mx} and renders the 4D$\to$3D unfolding of each
polytope as a \texttt{Graphics3D} object.
It depends on \texttt{unfold3DExport.m} for the core
SVD\,+\,Procrustes algorithm (Steps~1--3 above).

\paragraph{Fixed starting cell.}
By equivariance, all starting cells $C_1$ yield topologically
equivalent results.
The script therefore fixes $C_1 = 3$ (\texttt{fixedTop = 3}) and
enumerates every distinct successful peeling order
(deduplicated by order, not by $(C_1,C_2)$ pair).

\paragraph{Colour coding.}
Cells are coloured by position in the peeling order:
\begin{itemize}
  \item $C_1$: blue~(RGB $0.20, 0.45, 0.80$)
  \item middle cells: green~(RGB $0.15, 0.70, 0.40$)
  \item last cell: red~(RGB $0.90, 0.25, 0.15$)
\end{itemize}
Intermediate colours are produced by linear blending (\texttt{Blend}).

\paragraph{5-cell coordinate correction.}
The stored \texttt{localVers} in \texttt{ans4DGlobal.mx} were computed
from the original \texttt{f5.m}, which contained an incorrect fifth
vertex $v_5 \approx (0.191,\ldots)$ instead of the correct
$v_5 = \tfrac{1-\varphi}{2}(1,1,1,1) \approx (-0.309,\ldots)$.
The correct value satisfies $|v_k - v_5|^2 = 2$ for all $k=1,\ldots,4$,
matching the edge length of the regular 5-cell.
The file \texttt{f5.m} has been corrected (2026-05-08); the function
\texttt{recomputeLocalVers} recomputes the face-up--rotated coordinates
from the current vertex data so that the visualisation shows regular
tetrahedra even before \texttt{ans4DGlobal.mx} is regenerated.

\paragraph{Usage.}
\begin{verbatim}
SetDirectory["/path/to/260324Peeling4D"];
Get["unfold4D.m"]
\end{verbatim}
The script prints each polytope's nets sequentially in the notebook output.
The parameter \texttt{netImageSize} (default 500\,px) controls image size.


% ============================================================
\section{One-Step Backtrack Extension (\texttt{peeling1back.m})}
% ============================================================

The greedy algorithm fails when no candidate remains (neither left-side nor
fallback).
A \textbf{one-step backtrack} extension implemented in
\texttt{peeling1back.m} addresses this.

\subsection{Algorithm}

After the same preprocessing as the greedy version (centroid shift, rotation,
centroid computation, adjacency construction), the selection loop proceeds
as follows.

\begin{algorithm}[h]
\caption{peeling1back (one-step backtrack, selection loop)}
\label{alg:1back}
\begin{algorithmic}[1]
\Require Cell centroid array $\hat{c}$ (indexed $1\ldots|\mathit{cells}|$, each $\in\mathbb{R}^4$),
         adjacency list $\mathit{adj}$,
         starting cell $\mathit{top}$, second cell $C_2$
\Ensure  $(\mathit{order},\; \mathit{success})$
\State $\mathit{order} \gets [\mathit{top},\; C_2]$
\State $\mathit{nbrs} \gets$ copy of $\mathit{adj}$, with $\mathit{top}$ and $C_2$ removed
\State $\mathit{last} \gets C_2$
\State $\mathit{justBT} \gets \mathrm{False}$
  \Comment{$\mathit{justBT}$: True if the previous action was a backtrack}
\State $\mathit{tried} \gets \{\}$
  \Comment{$\mathit{tried}[\mathit{pos}]$: set of cells already tried at position $\mathit{pos}$}
\While{$|\mathit{order}| < |\mathit{cells}|$}
  \State $\mathit{pos} \gets |\mathit{order}|$
    \Comment{Current length of the selection sequence}
  \State $\mathit{cands} \gets \mathit{nbrs}[\mathit{last}]
           \setminus \mathit{tried}[\mathit{pos}]$
    \Comment{Untried candidates at the current position}
  \If{$\mathit{cands} = \emptyset$}
    \Comment{Dead end}
    \If{$\mathit{justBT}$ \textbf{or} $|\mathit{order}| \leq 2$}
      \Break
        \Comment{Do not backtrack more than one level}
    \EndIf
    \State $\mathit{bad} \gets \mathit{order}.\mathrm{last}()$
      \Comment{The cell that led to the dead end}
    \State $\mathit{order}.\mathrm{removeLast}()$
    \State $\mathit{nbrs} \gets$ restored to the state before the last forward step
    \State $\mathit{tried}[\mathit{pos}-1] \gets
             \mathit{tried}[\mathit{pos}-1] \cup \{\mathit{bad}\}$
      \Comment{Mark $\mathit{bad}$ as tried at the previous position}
    \State $\mathit{justBT} \gets \mathrm{True}$
  \Else
    \Comment{Advance}
    \State $L \gets \bigl\{j \in \mathit{cands} \mid
               \hat{c}_{\mathit{last},y}\,\hat{c}_{j,x}
               - \hat{c}_{\mathit{last},x}\,\hat{c}_{j,y} \geq 0\bigr\}$
    \If{$L \neq \emptyset$}
      \State $\mathit{next} \gets \arg\max_{j \in L}\; \hat{c}_{j,w}$
    \Else
      \State $\mathit{next} \gets \arg\min_{j \in \mathit{cands}}\; \hat{c}_{j,w}$
    \EndIf
    \State Save the current state of $\mathit{nbrs}$
    \State $\mathit{order}.\mathrm{append}(\mathit{next})$
    \State $\mathit{nbrs}[i] \gets \mathit{nbrs}[i] \setminus \{\mathit{next}\}$
      \hspace{1em}\textbf{for all} $i$
    \State $\mathit{last} \gets \mathit{next}$
    \State $\mathit{justBT} \gets \mathrm{False}$
  \EndIf
\EndWhile
\State \Return $\bigl(\mathit{order},\; |\mathit{order}| = |\mathit{cells}|\bigr)$
\end{algorithmic}
\end{algorithm}

This extension may succeed on some $(C_1, C_2)$ pairs where the greedy
version fails (quantitative comparison is a future task).


% ============================================================
\section{Future Work}
\label{sec:future}
% ============================================================

\begin{enumerate}
  \item \textbf{Complete the 600-cell computation.}
        The 600-cell is computationally expensive (600 cells;
        2,400 $(C_1,C_2)$ pairs in total);
        a complete classification has not yet been obtained.
        Parallelisation or additional computing resources are needed.

  \item \textbf{Tabulate overlap statistics.}
        For the 570 generated STL files, count the valid nets
        (no self-intersection) per polytope and summarise in a table.

  \item \textbf{Compare with the one-step backtrack extension.}
        Quantify the improvement in success rate achieved by
        \texttt{peeling1back.m} over the greedy \texttt{peeling4Df}.

  \item \textbf{Compare with Kaino (2019).}
        Verify whether the unfoldings produced by our algorithm are
        consistent with Kaino's geometrically constructed nets.

  \item \textbf{Symmetry of left-handed peeling.}
        Investigate the symmetry and duality between right-handed peeling
        ($c_y c'_x - c_x c'_y \geq 0$) and left-handed peeling
        ($c_y c'_x - c_x c'_y \leq 0$).

  \item \textbf{Systematic comparison of selection rules.}
        Apply Rules~1, 2, and~3 (Definition~\ref{defn:rules}) to small
        polytopes (5-cell, 8-cell, etc.) and compare success rates.
        Verification on 3D polyhedra (e.g.\ Archimedean solids) is also
        informative.

  \item \textbf{Implement the rigorous 4D left condition.}
        Replace the 2-component left condition
        $c_{k,y}c'_x - c_{k,x}c'_y \ge 0$ in the current
        \texttt{peeling4Df.m} with the determinant condition of
        Definition~\ref{defn:4dleft},
        $\det(\tilde{\mathbf{c}}_{k-1}, \tilde{\mathbf{c}}_k, \tilde{\mathbf{c}}_j) \ge 0$,
        and measure the change in success rate
        (paying attention to the $k=1$ initial condition).

  \item \textbf{Implement and evaluate two-stage selection.}
        Implement the two-stage strategy of Table~\ref{tab:twostage}
        (Rule~3 first, fall back to Rule~2) and compare its success rate
        with the greedy and one-step-backtrack versions.
\end{enumerate}


% ============================================================
\appendix
\section{Complete Code for \texttt{peeling4Df}}
% ============================================================

\begin{lstlisting}
peeling4Df[origVers_, faces_, cells_, top_] :=
  Module[
    {wHat, cc1, localVers, cCenters,
     neighbors, neighborsInit,
     order, cansInit, chosenFrom1,
     lastChosen, newChoice, cans,
     angle4D, rCans, cellVerts},

    wHat = {0, 0, 0, 1};
    (* wHat: unit vector along the +w axis (peeling axis) *)

    (* cellVerts[i]: set of vertex indices belonging to cell i *)
    cellVerts[i_] := Union @@ faces[[ cells[[i]] ]];

    (* Step 1: translate so that the centroid of all vertices is the origin *)
    localVers = (# - Mean[origVers]) & /@ origVers;

    (* Step 2: centroid of starting cell C1 in the translated frame *)
    cc1 = Mean[ localVers[[ cellVerts[top] ]] ];

    (* Step 3: rotate so that cc1 points in the +w direction *)
    localVers = Which[
      Abs[(cc1/Norm[cc1]).wHat - 1] < 0.0001, localVers,
      Abs[(cc1/Norm[cc1]).wHat + 1] < 0.0001, -localVers,
      True,
        angle4D = ArcCos[Clip[cc1.wHat/Norm[cc1], {-1,1}]];
        RotationMatrix[angle4D, {cc1/Norm[cc1], wHat}] . # & /@ localVers
    ];

    (* Step 4: centroids of all cells after rotation *)
    (* cCenters[[i]]: 4D centroid of cell i *)
    cCenters = Mean[localVers[[ cellVerts[#] ]]] & /@ Range[Length[cells]];

    (* Step 5: build adjacency list (cells sharing >= 1 face are adjacent) *)
    (* neighbors[[i]]: list of cell indices adjacent to cell i *)
    neighbors = Table[
      Select[Range[Length[cells]],
        Function[j, j!=i &&
          Length[Intersection[cells[[i]], cells[[j]]]] >= 1]],
      {i, Length[cells]}];

    (* Step 6: remove starting cell C1 from all adjacency lists *)
    neighbors = neighborsInit = Complement[#, {top}] & /@ neighbors;

    (* Step 7: C2 candidates = neighbours of C1 *)
    cansInit = neighbors[[top]];

    (* Step 8: peel for each C2 candidate *)
    Table[
      order = {top}; neighbors = neighborsInit;
      chosenFrom1 = cansInit[[i]];
      (* chosenFrom1: the chosen second cell C2 *)
      AppendTo[order, chosenFrom1];
      neighbors = Complement[#, {chosenFrom1}] & /@ neighbors;
      lastChosen = chosenFrom1; newChoice = lastChosen;

      While[
        Length[Flatten[neighbors]] != 0 &&
        Length[neighbors[[newChoice]]] != 0,

        (* Left half-space condition: c_y*c'_x - c_x*c'_y >= 0 *)
        (* cans: left half-space candidates *)
        cans = Select[neighbors[[lastChosen]],
          Function[j, With[{c = cCenters[[lastChosen]]},
            c[[2]]*cCenters[[j]][[1]] -
            c[[1]]*cCenters[[j]][[2]] >= 0]]];

        (* newChoice: next cell to be selected *)
        newChoice = Which[
          Length[cans] > 1,
            (* select left candidate with largest w-coordinate *)
            cans[[ First[First[
              Position[cCenters[[cans,4]], Max[cCenters[[cans,4]]]]
            ]] ]],
          Length[cans] == 1, cans[[1]],
          True,  (* fallback: right-side candidate with smallest w *)
            With[{rc = neighbors[[lastChosen]]},
              rc[[ First[First[
                Position[cCenters[[rc,4]], Min[cCenters[[rc,4]]]]
              ]] ]]]
        ];

        lastChosen = newChoice;
        AppendTo[order, newChoice];
        neighbors = Complement[#, {newChoice}] & /@ neighbors;
      ];

      {localVers, order, Length[order] == Length[cells]}
      , {i, Length[cansInit]}]
  ]
\end{lstlisting}


\begin{thebibliography}{9}

\bibitem{applePeeling}
  T.~Yoshino, S.~Chaidee, P.~Sawatdithep,
  \textit{Apple-Peel Unfolding of Polyhedra},
  \texttt{applePeelingV01.pdf}.

\bibitem{Kaino2019}
  K.~Kaino,
  \textit{Unfolding of Regular Polytopes},
  Symmetry Congress, Kanazawa (2019),
  \texttt{Kaino2019\_SymmetryCongressKanazawa.pdf}.

\end{thebibliography}

\end{document}
